%% ===========================================================================
\section{Experimental calibration protocol and limitations}
\label{sec:calibration}
%% ===========================================================================

The model contains four parameters that no constitutive table supplies: the
window size $a$, the fractional order $\alpha$, the relaxation time $\tau$ and
the flexoelectric coefficient $\mu$. This section states how each is to be
obtained, over which ranges the results of this paper hold, and which numbers
in the paper are not verified.

\subsection{Calibration of the window size and of the flexoelectric
coefficient}
\label{sec:calib_a_mu}

The window is not a fitting parameter of the discretisation: it is the second
moment of the averaging domain, and the domain is a physical object -- a
representative volume of the microstructure, or the range over which the
polarization is taken to be a continuum field. Two routes are available. The
geometric one takes $a$ from the micrograph: for a grain or fibre texture the
half-widths are the second moments of the grain-size distribution, which is
exactly the interpretation that a fabric tensor carries in
Sec.~\ref{sec:window}. The mechanical one takes $a$ from a size-effect
measurement on the same material: the frequency ratio of two otherwise
identical specimens of thickness $h_1$ and $h_2$ fixes the combination
$M_{22}/h^{2}$ through the boundary-layer solution of
Sec.~\ref{sec:bc_mech}, and the plate geometry then fixes $M_{22}$ itself.

The flexoelectric coefficient is measured, not assumed. The direct route is a
four-point bending test on a dielectric bar with electrodes on the two faces,
in which the open-circuit voltage under a known curvature $\kappa$ gives
$\mu$ through $V=\mu\kappa h^{2}/(\kappa_{yy}\varepsilon_0)$ for the geometry of
Fig.~\ref{fig:geometry}(a); the inverse route uses the same expression with the
measured voltage of a device whose curvature is known from its deflection. In
both routes the measured quantity is $\mu^{2}$ scaled data, so a single
measurement fixes the magnitude but not the sign, and the sign has to be taken
from the crystal symmetry. The value used in this paper is
$\mu=10^{-6}$\,\si{\coulomb\per\metre}, a nominal one; every bending coupling
in Secs.~\ref{sec:thickness} and~\ref{sec:transient} scales with $\mu^{2}$ and
should be read as a shape, not as a magnitude, until a measured coefficient is
substituted.

\subsection{Calibration of the fractional order and of the relaxation time}
\label{sec:calib_alpha_tau}

The pair $(\alpha,\tau)$ is obtained from a dielectric loss measurement. In the
frequency domain the fractional Kelvin--Voigt element \eqref{eq:kvfreq} has
$\tan\delta = (\omega\tau)^{\alpha}\sin(\pi\alpha/2)/
[1+(\omega\tau)^{\alpha}\cos(\pi\alpha/2)]$, so a log--log plot of
$\tan\delta(\omega)$ over two decades gives $\alpha$ as the low-frequency
slope and $\tau$ from the frequency at which $\tan\delta$ reaches its maximum.
The ring-down of a driven device gives the same information in the time domain:
the logarithmic decrement of \eqref{eq:zeta} is $\pi(\omega\tau)^{\alpha}
\sin(\pi\alpha/2)/2$ for small loss, so a measurement at two amplitudes fixes
the pair.

Two cautions follow from Sec.~\ref{sec:memorystudy}. First, the calibration
must state which matrix the memory multiplies, because the same measured
decrement corresponds to two different $(\alpha,\tau)$ pairs under
$K_{\rm eff}$ and under $K_{uu}$ -- the ratio of the resulting damping ratios
is $0.449$ at $\alpha=1$. Second, the fit is valid where the measurement is:
the fractional element describes a broadened, single-power-law loss, and a
material with two well-separated loss peaks needs two elements rather than one.

\subsection{Validity window}
\label{sec:limit_window}

Three inequalities bound the results, and each is a modelling assumption rather
than a numerical tolerance.

\begin{enumerate}
\item \emph{Small coupling at the operator level.} The mode-free coupling
  factor \eqref{eq:kstatic} is used with the uncoupled denominator whenever
  $k^{2}\ll1$; for the extensional configuration of this paper
  $k^{2}=0.5581$, which is \emph{not} small, and the coupled denominator is
  used instead. The uncoupled form is retained only for the comparison with the
  constitutive closed form, where it is stated explicitly.
\item \emph{Small window.} The second-order truncation \eqref{eq:Mdef} of the
  averaging operator is valid while the window is small against the length over
  which the strain gradient varies, which for the plate of
  Table~\ref{tab:params} means $a/L\le0.2$. The bound is set by
  Fig.~\ref{fig:domain}: beyond it the retained second moment no longer
  represents the average, independently of the mesh.
\item \emph{Operating band.} The damping map is claimed for
  $\omega\tau\le1$, the band in which the spring-pot represents a relaxation
  whose spectrum is centred on the device frequency. Beyond that band the map
  is an extrapolation of the fitted law, not a prediction, and
  Table~\ref{tab:modulus} shows how far the integer-order alternative departs
  from it there.
\end{enumerate}

\subsection{Limitations}
\label{sec:limitations}

\emph{Two dimensions.} The formulation is two-dimensional, with a plane-stress
reduction of the constitutive tensors. The reduction is a modelling choice:
plane strain would change the extensional coupling, and the three-dimensional
solution differs from both. The thin-plate results of
Sec.~\ref{sec:thickness} are therefore two-dimensional results, consistent with
the plane-stress constitutive law that the table states.

\emph{Unverified parameters.} Two inputs are nominal rather than measured. The
mass density $\rho=\SI{6020}{\kilo\gram\per\cubic\metre}$ scales every
frequency in the paper by a common factor; no conclusion depends on its value,
but the absolute kilohertz numbers in Secs.~\ref{sec:orientation}
and~\ref{sec:application} should be rescaled if a measured density is
substituted. The flexoelectric coefficient $\mu=10^{-6}$\,\si{\coulomb\per\metre}
scales every flexoelectric coupling as $\mu^{2}$; the growth factor, the
saturation and the comparison with the piezoelectric response are unaffected by
its value, but the absolute coupling is not.

\emph{Numbers set by the solver.} The two-significant-digit statement of
Sec.~\ref{sec:thickness} is a limitation as much as a result: the third digit
of the thin-plate coupling is not a property of the model. The same applies,
with smaller amplitude, to any quantity in the paper obtained as a difference
of nearly equal eigenvalues.

\emph{Scope of the validation.} The verification of
Sec.~\ref{sec:verification} and Appendix~\ref{app:verification} establishes
that the implementation solves the stated equations; it does not establish that
the equations describe a particular material. That step requires the
measurements of Secs.~\ref{sec:calib_a_mu} and~\ref{sec:calib_alpha_tau}, and
it is the natural continuation of this work.
